\documentclass[aspectratio=169]{beamer}
\usepackage[utf8]{inputenc}
\usepackage[T1]{fontenc}
\usepackage[es-ES]{babel}
\usepackage{graphicx}
\usepackage{tikz}
\usepackage{listings}
\usepackage{xcolor}

% Tema
\usetheme{Madrid}
\usecolortheme{default}

% Colores personalizados
\definecolor{unsazul}{RGB}{30, 60, 120}
\definecolor{unsverde}{RGB}{0, 100, 0}
\definecolor{unsrojo}{RGB}{200, 100, 100}

\setbeamercolor{structure}{fg=unsazul}
\setbeamercolor{palette primary}{bg=unsazul,fg=white}
\setbeamercolor{palette secondary}{bg=unsazul!80,fg=white}
\setbeamercolor{palette tertiary}{bg=unsazul!60,fg=white}

% Configuración de listings
\lstset{
    basicstyle=\ttfamily\small,
    keywordstyle=\color{unsazul}\bfseries,
    commentstyle=\color{gray}\itshape,
    stringstyle=\color{unsverde},
    showstringspaces=false,
    breaklines=true,
    frame=single,
    backgroundcolor=\color{gray!10}
}

% Información del documento
\title{Métodos de Análisis de Datos 1}
\subtitle{Clase 7: APIs, Web Scraping y Tidy Data}
\author{Universidad Nacional del Sur}
\date{}

\begin{document}

% Portada
\begin{frame}
\titlepage
\end{frame}

% Contenidos
\begin{frame}{Contenidos de hoy}
\tableofcontents
\end{frame}

\section{APIs REST}

\begin{frame}{¿Qué es una API REST?}
\begin{block}{API: Application Programming Interface}
Interfaz que permite comunicación entre aplicaciones
\end{block}

\begin{block}{REST: REpresentational State Transfer}
Estilo de arquitectura para servicios web
\end{block}

\pause

\begin{exampleblock}{Ejemplos de APIs públicas}
\begin{itemize}
    \item OpenWeatherMap: datos meteorológicos
    \item Twitter/X: tweets, usuarios
    \item GitHub: repositorios, commits
    \item Google Maps: mapas, geocodificación
    \item The Movie Database: películas, series
\end{itemize}
\end{exampleblock}
\end{frame}

\begin{frame}{Componentes de una API REST}
\begin{enumerate}
    \item \textbf{Endpoint (URL):} Dirección del recurso
    \begin{itemize}
        \item \texttt{https://api.ejemplo.com/users}
    \end{itemize}
    
    \item \textbf{Métodos HTTP:}
    \begin{itemize}
        \item GET: obtener datos
        \item POST: crear nuevos datos
        \item PUT/PATCH: actualizar datos
        \item DELETE: eliminar datos
    \end{itemize}
    
    \item \textbf{Headers:} Metadatos (autenticación, tipo de contenido)
    
    \item \textbf{Parámetros:} Query string o body
    
    \item \textbf{Respuesta:} Generalmente JSON o XML
\end{enumerate}
\end{frame}

\begin{frame}{Anatomía de una request HTTP}
\begin{exampleblock}{Ejemplo completo}
\texttt{GET https://api.openweathermap.org/data/2.5/weather?q=BahiaBlanca\&appid=ABC123}
\end{exampleblock}

\begin{itemize}
    \item \textbf{Método:} GET
    \item \textbf{Base URL:} \texttt{https://api.openweathermap.org}
    \item \textbf{Endpoint:} \texttt{/data/2.5/weather}
    \item \textbf{Query params:}
    \begin{itemize}
        \item \texttt{q=BahiaBlanca} (ciudad)
        \item \texttt{appid=ABC123} (API key)
    \end{itemize}
\end{itemize}

\pause

\begin{block}{Respuesta típica (JSON)}
\texttt{\{"temp": 25.3, "humidity": 60, "description": "clear sky"\}}
\end{block}
\end{frame}

\begin{frame}{Códigos de respuesta HTTP}
\begin{columns}[t]
\column{0.5\textwidth}
\textbf{Éxito (2xx):}
\begin{itemize}
    \item \textcolor{unsverde}{\textbf{200 OK}}: Exitoso
    \item \textcolor{unsverde}{\textbf{201 Created}}: Recurso creado
    \item \textcolor{unsverde}{\textbf{204 No Content}}: Sin contenido
\end{itemize}

\textbf{Redirección (3xx):}
\begin{itemize}
    \item \textbf{301 Moved}: Movido permanentemente
\end{itemize}

\column{0.5\textwidth}
\textbf{Error del cliente (4xx):}
\begin{itemize}
    \item \textcolor{unsrojo}{\textbf{400 Bad Request}}: Solicitud inválida
    \item \textcolor{unsrojo}{\textbf{401 Unauthorized}}: No autenticado
    \item \textcolor{unsrojo}{\textbf{404 Not Found}}: No encontrado
    \item \textcolor{unsrojo}{\textbf{429 Too Many}}: Rate limit
\end{itemize}

\textbf{Error del servidor (5xx):}
\begin{itemize}
    \item \textcolor{unsrojo}{\textbf{500 Internal Error}}: Error del servidor
\end{itemize}
\end{columns}
\end{frame}

\begin{frame}[fragile]{GET Request en Python}
\begin{lstlisting}[language=Python]
import requests
import pandas as pd

# Request simple
url = "https://api.ejemplo.com/posts"
response = requests.get(url)

# Verificar éxito
if response.status_code == 200:
    data = response.json()
    df = pd.DataFrame(data)
    print(df.head())
else:
    print(f"Error: {response.status_code}")

# Con parámetros
params = {'userId': 1, 'limit': 10}
response = requests.get(url, params=params)

# Con autenticación
headers = {'Authorization': 'Bearer TU_TOKEN'}
response = requests.get(url, headers=headers)
\end{lstlisting}
\end{frame}

\begin{frame}[fragile]{GET Request en R}
\begin{lstlisting}[language=R]
library(httr)
library(jsonlite)

# Request simple
url <- "https://api.ejemplo.com/posts"
response <- GET(url)

# Verificar éxito
if (status_code(response) == 200) {
  data <- content(response, as = "text")
  df <- fromJSON(data)
  head(df)
} else {
  print(paste("Error:", status_code(response)))
}

# Con parámetros
response <- GET(url, query = list(userId = 1, limit = 10))

# Con autenticación
response <- GET(url, add_headers(Authorization = "Bearer TU_TOKEN"))
\end{lstlisting}
\end{frame}

\begin{frame}[fragile]{POST Request (crear datos)}
\begin{columns}[t]
\column{0.5\textwidth}
\textbf{Python:}
\begin{lstlisting}[language=Python]
import requests

url = "https://api.ejemplo.com/posts"

# Datos a enviar
nuevo_post = {
    'title': 'Nuevo post',
    'body': 'Contenido',
    'userId': 1
}

# POST request
response = requests.post(
    url, 
    json=nuevo_post
)

if response.status_code == 201:
    print("Creado:", response.json())
\end{lstlisting}

\column{0.5\textwidth}
\textbf{R:}
\begin{lstlisting}[language=R]
library(httr)

url <- "https://api.ejemplo.com/posts"

# Datos a enviar
nuevo_post <- list(
  title = 'Nuevo post',
  body = 'Contenido',
  userId = 1
)

# POST request
response <- POST(
  url, 
  body = nuevo_post, 
  encode = "json"
)

if (status_code(response) == 201) {
  print("Creado")
}
\end{lstlisting}
\end{columns}
\end{frame}

\begin{frame}{Rate Limiting}
\begin{alertblock}{Problema}
Las APIs limitan el número de requests por tiempo
\end{alertblock}

\begin{exampleblock}{Límites típicos}
\begin{itemize}
    \item 100 requests por hora
    \item 1000 requests por día
    \item 10 requests por segundo
\end{itemize}
\end{exampleblock}

\pause

\begin{block}{Soluciones}
\begin{enumerate}
    \item Espaciar requests (sleep entre llamadas)
    \item Cachear resultados
    \item Pagar por plan premium
    \item Usar múltiples API keys (con cuidado)
\end{enumerate}
\end{block}
\end{frame}

\begin{frame}[fragile]{Implementar Rate Limiting}
\begin{columns}[t]
\column{0.5\textwidth}
\textbf{Python:}
\begin{lstlisting}[language=Python]
import time
import requests

urls = [
    'url1', 'url2', 'url3'
]

resultados = []
for url in urls:
    response = requests.get(url)
    resultados.append(response.json())
    
    # Esperar 1 segundo
    time.sleep(1)
\end{lstlisting}

\column{0.5\textwidth}
\textbf{R:}
\begin{lstlisting}[language=R]
library(httr)

urls <- c('url1', 'url2', 'url3')

resultados <- list()
for (url in urls) {
  response <- GET(url)
  resultados[[length(resultados) + 1]] <- content(response)
  
  # Esperar 1 segundo
  Sys.sleep(1)
}
\end{lstlisting}
\end{columns}
\end{frame}

\section{Web Scraping}

\begin{frame}{¿Qué es web scraping?}
\begin{block}{Definición}
Extraer datos de páginas web mediante programación
\end{block}

\begin{alertblock}{¿Cuándo usar?}
\textbf{SOLO cuando no hay API disponible}
\end{alertblock}

\pause

\begin{block}{Proceso}
\begin{enumerate}
    \item Descargar HTML de la página
    \item Inspeccionar estructura (DevTools del navegador)
    \item Identificar selectores CSS o XPath
    \item Extraer información específica
    \item Limpiar y estructurar datos
\end{enumerate}
\end{block}
\end{frame}

\begin{frame}{Herramientas de web scraping}
\begin{columns}[t]
\column{0.5\textwidth}
\textbf{Python:}
\begin{itemize}
    \item \textbf{BeautifulSoup}: Parsing HTML simple
    \item \textbf{Scrapy}: Framework completo
    \item \textbf{Selenium}: Páginas dinámicas (JavaScript)
    \item \textbf{requests-html}: Combina requests y BS4
\end{itemize}

\column{0.5\textwidth}
\textbf{R:}
\begin{itemize}
    \item \textbf{rvest}: Principal librería
    \item \textbf{xml2}: Parsing XML/HTML
    \item \textbf{RSelenium}: Páginas dinámicas
\end{itemize}
\end{columns}

\vspace{1em}

\begin{exampleblock}{Recomendación}
Comenzar con BeautifulSoup (Python) o rvest (R)
\end{exampleblock}
\end{frame}

\begin{frame}{Selectores CSS básicos}
\begin{block}{¿Qué son?}
Patrones para identificar elementos HTML
\end{block}

\begin{center}
\begin{tabular}{|l|l|}
\hline
\textbf{Selector} & \textbf{Descripción} \\
\hline
\texttt{p} & Todos los párrafos \\
\texttt{.clase} & Elementos con clase="clase" \\
\texttt{\#id} & Elemento con id="id" \\
\texttt{div p} & Párrafos dentro de divs \\
\texttt{div > p} & Párrafos hijos directos de divs \\
\texttt{a[href]} & Enlaces con atributo href \\
\texttt{.precio.destacado} & Elementos con ambas clases \\
\hline
\end{tabular}
\end{center}

\vspace{0.5em}
\small{Usar DevTools del navegador (F12) para inspeccionar elementos}
\end{frame}

\begin{frame}[fragile]{Web Scraping en Python}
\begin{lstlisting}[language=Python]
import requests
from bs4 import BeautifulSoup

# Descargar página
url = "https://ejemplo.com/productos"
response = requests.get(url)
soup = BeautifulSoup(response.content, 'html.parser')

# Extraer todos los enlaces
enlaces = soup.find_all('a')
for link in enlaces:
    print(link.get('href'))

# Extraer tabla
tabla = soup.find('table', {'class': 'datos'})
filas = tabla.find_all('tr')

# Extraer precios por clase
precios = soup.find_all('span', {'class': 'precio'})
for precio in precios:
    print(precio.text.strip())

# Con selector CSS
productos = soup.select('.producto .nombre')
\end{lstlisting}
\end{frame}

\begin{frame}[fragile]{Web Scraping en R}
\begin{lstlisting}[language=R]
library(rvest)

# Descargar página
url <- "https://ejemplo.com/productos"
pagina <- read_html(url)

# Extraer enlaces
enlaces <- pagina %>%
  html_nodes("a") %>%
  html_attr("href")

# Extraer tabla
tabla <- pagina %>%
  html_node("table.datos") %>%
  html_table()

df <- as_tibble(tabla)

# Extraer precios
precios <- pagina %>%
  html_nodes(".precio") %>%
  html_text()
\end{lstlisting}
\end{frame}

\begin{frame}{Consideraciones legales y éticas}
\begin{alertblock}{SIEMPRE verificar antes de scrapear}
\begin{enumerate}
    \item \textbf{Términos de servicio:} ¿Permiten scraping?
    \item \textbf{robots.txt:} \texttt{https://sitio.com/robots.txt}
    \item \textbf{Frecuencia:} No saturar el servidor
    \item \textbf{User-Agent:} Identificarse apropiadamente
    \item \textbf{Datos personales:} Respetar privacidad
\end{enumerate}
\end{alertblock}

\pause

\begin{exampleblock}{Buenas prácticas}
\begin{itemize}
    \item Pausas entre requests (1-3 segundos)
    \item Cachear resultados
    \item Scrapear de noche (menos carga)
    \item Respetar directivas de robots.txt
    \item Contactar al sitio si tenés dudas
\end{itemize}
\end{exampleblock}
\end{frame}

\begin{frame}[fragile]{User-Agent y robots.txt}
\textbf{Configurar User-Agent:}
\begin{lstlisting}[language=Python]
# Python
headers = {
    'User-Agent': 'MiScraper/1.0 (contacto@ejemplo.com)'
}
response = requests.get(url, headers=headers)
\end{lstlisting}

\vspace{0.5em}

\textbf{Verificar robots.txt:}
\begin{verbatim}
https://wikipedia.org/robots.txt

User-agent: *
Disallow: /wiki/Special:
Allow: /wiki/
\end{verbatim}

\begin{alertblock}{Importante}
Si robots.txt prohíbe, \textbf{NO scrapear}
\end{alertblock}
\end{frame}

\section{Tidy Data}

\begin{frame}{Principios de Tidy Data}
\begin{block}{¿Qué son datos ordenados (tidy)?}
Datos estructurados de forma consistente para facilitar análisis
\end{block}

\begin{block}{Tres principios (Hadley Wickham)}
\begin{enumerate}
    \item \textbf{Cada variable forma una columna}
    \item \textbf{Cada observación forma una fila}
    \item \textbf{Cada tipo de unidad observacional forma una tabla}
\end{enumerate}
\end{block}

\pause

\begin{alertblock}{¿Por qué importa?}
\begin{itemize}
    \item Facilita manipulación con dplyr/pandas
    \item Visualización más simple
    \item Menos errores
    \item Código más claro y mantenible
\end{itemize}
\end{alertblock}
\end{frame}

\begin{frame}{Ejemplo: Datos NO ordenados}
\textbf{Wide format (ancho):}

\begin{center}
\begin{tabular}{|l|c|c|c|}
\hline
\textbf{Persona} & \textbf{2020} & \textbf{2021} & \textbf{2022} \\
\hline
Ana & 50000 & 52000 & 54000 \\
Luis & 48000 & 49000 & 51000 \\
\hline
\end{tabular}
\end{center}

\vspace{0.5em}

\begin{alertblock}{Problemas}
\begin{itemize}
    \item Los años están en columnas (deberían ser valores)
    \item Difícil agregar nuevos años
    \item Difícil visualizar series temporales
\end{itemize}
\end{alertblock}
\end{frame}

\begin{frame}{Ejemplo: Datos ordenados}
\textbf{Long format (largo):}

\begin{center}
\begin{tabular}{|l|c|c|}
\hline
\textbf{Persona} & \textbf{Año} & \textbf{Ingreso} \\
\hline
Ana & 2020 & 50000 \\
Ana & 2021 & 52000 \\
Ana & 2022 & 54000 \\
Luis & 2020 & 48000 \\
Luis & 2021 & 49000 \\
Luis & 2022 & 51000 \\
\hline
\end{tabular}
\end{center}

\vspace{0.5em}

\begin{block}{Ventajas}
\begin{itemize}
    \item Fácil filtrar por año
    \item Fácil agrupar y sumarizar
    \item Compatible con ggplot2/seaborn
\end{itemize}
\end{block}
\end{frame}

\section{Reshape: Pivot y Melt}

\begin{frame}{Wide → Long (Pivot/Melt)}
\begin{block}{Objetivo}
Convertir columnas en filas (wide → long)
\end{block}

\begin{center}
\begin{tikzpicture}[scale=0.8]
    % Wide
    \node at (0,3) {\textbf{Wide}};
    \draw (0,0) rectangle (4,2);
    \node at (2,1.5) {Persona | 2020 | 2021 | 2022};
    \node at (2,1) {Ana | 50k | 52k | 54k};
    \node at (2,0.5) {Luis | 48k | 49k | 51k};
    
    % Arrow
    \draw[->, thick] (4.5,1) -- (5.5,1);
    
    % Long
    \node at (8,3) {\textbf{Long}};
    \draw (6,0) rectangle (10,2);
    \node at (8,1.5) {Persona | Año | Ingreso};
    \node at (8,1) {Ana | 2020 | 50k};
    \node at (8,0.5) {Ana | 2021 | 52k};
\end{tikzpicture}
\end{center}

\vspace{0.5em}

\begin{itemize}
    \item Python: \texttt{pd.melt()}
    \item R: \texttt{pivot\_longer()}
\end{itemize}
\end{frame}

\begin{frame}[fragile]{Wide → Long: Python}
\begin{lstlisting}[language=Python]
import pandas as pd

# Datos wide
df_wide = pd.DataFrame({
    'persona': ['Ana', 'Luis'],
    '2020': [50000, 48000],
    '2021': [52000, 49000],
    '2022': [54000, 51000]
})

# Melt (wide → long)
df_long = pd.melt(
    df_wide,
    id_vars=['persona'],           # Columnas a mantener
    value_vars=['2020', '2021', '2022'],  # Columnas a convertir
    var_name='año',                # Nombre de nueva columna
    value_name='ingreso'           # Nombre de valores
)

# Simplificado (todas excepto id_vars)
df_long = pd.melt(
    df_wide, 
    id_vars=['persona'], 
    var_name='año', 
    value_name='ingreso'
)
\end{lstlisting}
\end{frame}

\begin{frame}[fragile]{Wide → Long: R}
\begin{lstlisting}[language=R]
library(tidyverse)

# Datos wide
df_wide <- tibble(
  persona = c('Ana', 'Luis'),
  `2020` = c(50000, 48000),
  `2021` = c(52000, 49000),
  `2022` = c(54000, 51000)
)

# Pivot longer (wide → long)
df_long <- df_wide %>%
  pivot_longer(
    cols = c(`2020`, `2021`, `2022`),  # Columnas a convertir
    names_to = "año",                   # Nombre de columna
    values_to = "ingreso"               # Nombre de valores
  )

# Con patrón
df_long <- df_wide %>%
  pivot_longer(
    cols = matches("^\\d{4}$"),  # Regex para años
    names_to = "año",
    values_to = "ingreso"
  )
\end{lstlisting}
\end{frame}

\begin{frame}{Long → Wide (Pivot)}
\begin{block}{Objetivo}
Convertir filas en columnas (long → wide)
\end{block}

\begin{center}
\begin{tikzpicture}[scale=0.8]
    % Long
    \node at (2,3) {\textbf{Long}};
    \draw (0,0) rectangle (4,2);
    \node at (2,1.5) {Persona | Año | Ingreso};
    \node at (2,1) {Ana | 2020 | 50k};
    \node at (2,0.5) {Ana | 2021 | 52k};
    
    % Arrow
    \draw[->, thick] (4.5,1) -- (5.5,1);
    
    % Wide
    \node at (8,3) {\textbf{Wide}};
    \draw (6,0) rectangle (10,2);
    \node at (8,1.5) {Persona | 2020 | 2021 | 2022};
    \node at (8,1) {Ana | 50k | 52k | 54k};
    \node at (8,0.5) {Luis | 48k | 49k | 51k};
\end{tikzpicture}
\end{center}

\vspace{0.5em}

\begin{itemize}
    \item Python: \texttt{pivot()} o \texttt{pivot\_table()}
    \item R: \texttt{pivot\_wider()}
\end{itemize}
\end{frame}

\begin{frame}[fragile]{Long → Wide: Python}
\begin{lstlisting}[language=Python]
# Pivot (long → wide)
df_wide = df_long.pivot(
    index='persona',      # Columna para filas
    columns='año',        # Columna para columnas
    values='ingreso'      # Valores a rellenar
)

# Si hay duplicados, usar pivot_table con agregación
df_wide = df_long.pivot_table(
    index='persona',
    columns='año',
    values='ingreso',
    aggfunc='mean'  # Función de agregación
)

# Reset index para tener 'persona' como columna
df_wide = df_wide.reset_index()
\end{lstlisting}
\end{frame}

\begin{frame}[fragile]{Long → Wide: R}
\begin{lstlisting}[language=R]
# Pivot wider (long → wide)
df_wide <- df_long %>%
  pivot_wider(
    names_from = año,      # Columna que se convierte en columnas
    values_from = ingreso  # Valores a rellenar
  )

# Con prefijo
df_wide <- df_long %>%
  pivot_wider(
    names_from = año,
    values_from = ingreso,
    names_prefix = "año_"
  )

# Con múltiples columnas de valores
df_wide <- df_long %>%
  pivot_wider(
    names_from = año,
    values_from = c(ingreso, bonus)
  )
\end{lstlisting}
\end{frame}

\section{Resumen}

\begin{frame}{Resumen de la clase}
\begin{itemize}
    \item Las \textbf{APIs REST} proveen acceso estructurado a datos
    \item Componentes: endpoint, método HTTP, headers, parámetros, respuesta
    \item Códigos: 200 (OK), 401 (no autorizado), 404 (no encontrado)
    \item Implementar \textbf{rate limiting} para respetar límites
    \item \textbf{Web scraping}: último recurso cuando no hay API
    \item Herramientas: BeautifulSoup (Python), rvest (R)
    \item Respetar \textbf{robots.txt} y términos de servicio
    \item \textbf{Tidy data}: cada variable una columna, cada observación una fila
    \item \textbf{Wide → Long}: melt/pivot\_longer
    \item \textbf{Long → Wide}: pivot/pivot\_wider
\end{itemize}

\vspace{1em}

\begin{block}{Próxima unidad}
\textbf{Unidad 2: Exploración y Visualización de Datos}
\end{block}
\end{frame}

\begin{frame}{Tarea}
\begin{enumerate}
    \item \textbf{APIs:}
    \begin{itemize}
        \item Elegir una API pública (ej: JSONPlaceholder, OpenWeather)
        \item Hacer al menos 3 requests diferentes
        \item Convertir respuestas a DataFrames
    \end{itemize}
    
    \item \textbf{Web scraping (opcional):}
    \begin{itemize}
        \item Elegir un sitio simple (verificar robots.txt)
        \item Extraer una tabla o lista
        \item Estructurar en DataFrame
    \end{itemize}
    
    \item \textbf{Tidy data:}
    \begin{itemize}
        \item Crear dataset wide
        \item Convertir a long
        \item Volver a wide
        \item Visualizar ambos formatos
    \end{itemize}
    
    \item Documentar todo y subir a GitHub
\end{enumerate}
\end{frame}

\begin{frame}[standout]
\Huge ¿Preguntas?
\end{frame}

\end{document}
